\documentclass[11pt]{article}
\usepackage[margin=1in]{geometry}
\usepackage{amsmath,amssymb,mathtools}
\usepackage{booktabs,array,tabularx}
\usepackage[colorlinks=true,urlcolor=blue]{hyperref}
\newcommand{\Pj}{p^{\theta}_{j}}
\newcommand{\Ps}{p^{\ast}_{j}}
\newcommand{\ttil}{\tilde t}
\setlength{\parindent}{0pt}\setlength{\parskip}{5pt}
\begin{document}

\begin{center}
{\Large\bfseries What Pearson training identifies about a judge's rating distribution}\\[4pt]
11-21-2026
\end{center}

\textbf{Date:} 11-21-2026. \textbf{By:} ultimate\_brando9 (Brando's AI assistant), Claude Sonnet 5.5, Claude Code headless harness; checks run by sympy 1.14, numpy/scipy, and Lean 4.22 with Mathlib (\texttt{math/}). Codex added the first-moment explanation and the small-data study on 11-21-2026. This file is the typeset form of \texttt{MATH\_11-21-2026.md} (same math, same numbers).

\section*{Only the first conditional moment enters Pearson}
On each artifact $x$, a judge predicts six probabilities $q_x(r)$ for ratings $r=0,\ldots,5$.
Pearson correlation of mean ratings reads just $m_q(x)=\sum_r r q_x(r)$, the first conditional
moment. The map from probabilities to $m_q$ is linear. A judge can use nonlinear input features
to predict $m_q(x)$, and Pearson's normalization is nonlinear.
Even perfect correlation leaves a common offset and positive rescaling of the predicted means undetermined.

\textbf{What can go missing?} Uniform forecasts on $\{1,2,3\}$ and on $\{0,2,4\}$ both
have mean 2, but their variances are $2/3$ and $8/3$. Replacing one with the other on an artifact leaves
every mean prediction and hence every defined mean-based correlation unchanged.
With six categories, fixing total mass and mean leaves four independent directions in which the probabilities can vary at an interior distribution.

\textbf{What cross-entropy learns.} For a target distribution $p_x$, expected cross-entropy (CE)
has excess risk $D_{\mathrm{KL}}(p_x\Vert q_x)$ in nats. Its unique population optimum over all
distributions is $q_x=p_x$, so it recovers the mean and every higher moment. For any
integer $j\ge 1$, let $u_j(r)=(r/5)^j$. Since $0\le u_j\le1$, total variation distance (TVD)
and \href{https://www.stat.yale.edu/\string~yw562/teaching/598/lec04.pdf}{Pinsker’s inequality (Thm. 4.4)} give
\[
\left|\mathbb E_{p_x}u_j(R)-\mathbb E_{q_x}u_j(R)\right|
\le \mathrm{TVD}(p_x,q_x)\le\sqrt{D_{\mathrm{KL}}(p_x\Vert q_x)/2}.
\]
This bound uses population excess CE. With few training ratings, sampling noise and model
restrictions can change which fitted judge has lower held-out loss.
The \href{https://github.com/brando90/judges-dont-need-psychometrics-or-corr/tree/main/experiments/19_small_data_moment_recovery_open_and_closed_judges_start_21Nov2026}{small-data moment study} tests this separately for open-weight and closed-weight judges.

\section*{What correlation can and cannot learn (11-21-2026)}
\textbf{In one sentence.} The correlation objective reads exactly one number per item from the judge, the first moment (the mean) $m_i=\sum_r r\,p_i(r)$, and only up to a shift and a positive rescaling of the means. The spread, the shape and every higher moment of $p_i$ are invisible to it. CE, RPS and TVD read the whole distribution.

\textbf{1. The dependence is through the means only (R1).} Since $\sum_i\ttil_i=0$, the numerator of the box is $\sum_i\ttil_i m_i$, so
\[
\rho=\frac{\sum_i\ttil_i\,m_i}{\sqrt{\sum_i\ttil_i^{\,2}\cdot\sum_i(m_i-\bar m)^2}}=F(m_1,\dots,m_n),\qquad m_i=\sum_{r\in R} r\,p_i(r),
\]
a function of the $n$ numbers $m_i$ and of the fixed target means $\ttil_i$. The numerator is linear in the judge's distribution: the judge enters through the $n$ linear functionals $p_i\mapsto\sum_r r\,p_i(r)$.

\textbf{2. Level and scale are not identified (R2).} For $b>0$ and any $a$,
\[
\rho(a+b\,m)=\rho(m),\qquad\text{so}\qquad \rho=1\ \Longleftrightarrow\ m_i=a+b\,t_i\ \ (b>0).
\]
A judge whose means are $a+b\,t_i$ for any $a$ and any $b>0$ is correlation-optimal, whatever its level, scale or spread.

\textbf{3. Spread, shape and higher moments are invisible (R2).} The $k$-th moment $\mu_k(p_i)=\sum_r r^k p_i(r)$ and the variance $\sigma_i^2=\sum_r(r-m_i)^2p_i(r)=\mu_2(p_i)-m_i^2$ do not appear in $F$. Two judges with the same means $m_i$ get the same $\rho$ for every choice of $\sigma_i^2,\mu_3,\mu_4,\dots$ The pair above (uniform on $\{1,2,3\}$ and on $\{0,2,4\}$, Lean-checked, see the checks table) has the same mean $2$, hence the same $\rho$, while the variances are $2/3$ and $8/3$ and the RPS between them is $2/9>0$. Within a fixed mean, each item still has $K-2$ free shape dimensions, so the correlation-optimal set has $n(K-2)$ free dimensions plus $(a,b)$.

\textbf{4. The gradient lives on one direction per item (R3).} With $p_i=\mathrm{softmax}(u_i)$,
\[
\begin{aligned}
\frac{\partial\rho}{\partial u_i(r)}&=p_i(r)\,(r-m_i)\,g_i &&(\text{one fixed direction in }r,\ \text{scaled by the number }g_i),\\
\frac{\partial L_{CE}}{\partial u_i(r)}&=p_i(r)-p^\ast_i(r) &&(\text{all }K-1\text{ directions}).
\end{aligned}
\]
Correlation training can only push each item's mean; CE training can also correct its shape.

\textbf{5. What CE, RPS and TVD read (R4--R6).} CE and RPS are strictly proper and TVD vanishes only at the target: each is minimized only at $p_i=p^\ast_i$, and $p_i=p^\ast_i$ fixes the mean and every higher moment, $\mu_k(p_i)=\mu_k(p^\ast_i)$ for all $k$. CE splits exactly into noise, a shape gap $G(P)\ge0$ out of reach of every mean-only objective, and a calibration term (R4).

\begin{center}
\renewcommand{\arraystretch}{1.25}
\begin{tabular}{@{}lcccc@{}}
\toprule
 & mean $m_i$ & level, scale & variance, shape, $\mu_{k\ge3}$ & minimized only at $p^\ast$\\
\midrule
Corr ($\rho$) & read & not identified & not read & no (R2, R6)\\
CE & read & read & read & yes\\
RPS & read & read & read & yes\\
TVD & read & read & read & yes\\
\bottomrule
\end{tabular}
\end{center}

\textbf{What this shows, and what it does not.} It shows, at the population level (the human distribution $p^\ast_i$ known exactly, i.e. unlimited data), what each objective can see and where each is optimal. It does \emph{not} show the limited-data claim: whether a CE-trained judge also does better with few ratings is empirical (experiments 02, 08, 11, 16, 17 and the sweep in this folder), and a separate new experiment tests it.

\section*{Pearson between mean ratings}
Replace the placeholder on the page by
\[
\boxed{\ \mathrm{Corr}\big(\Pj,\Ps\big)=\frac{\sum_{i}\ttil_i\,\sum_{r\in R} r\,\Pj(r\mid x_i,y^\ast_i,\hat y_i;\theta)}{\sqrt{\sum_i \ttil_i^{\,2}\ \cdot\ \sum_i\big(m_i-\bar m\big)^2}}\ }
\]
The terms in this equation are
\begin{align*}
m_i&=\sum_{r\in R} r\,\Pj(r\mid x_i,y^\ast_i,\hat y_i;\theta) &&\text{the judge's mean rating on item } i,\\
t_i&=\sum_{r\in R} r\,\Ps(r\mid x_i,y^\ast_i,\hat y_i;h) &&\text{the target mean rating (human or persona } h),\\
\ttil_i&=t_i-\bar t,\quad \bar t=\tfrac1n\textstyle\sum_i t_i &&\text{the centered target mean},\\
\bar m&=\tfrac1n\textstyle\sum_i m_i &&\text{the average judge mean},\\
\textstyle\sum_i&=\textstyle\sum_{i=1}^{n} &&\text{the sum over items } i \ (\text{the page's } v\in V).
\end{align*}
The judge's distribution $\Pj$ enters only through the $n$ numbers $m_i$ (its mean ratings), in the numerator linearly (coefficient $r\,\ttil_i$ on $p_i(r)$) and in the denominator through their spread. Everything about the shape of $p_i$ beyond its mean is invisible. CE, RPS and TVD are functions of the whole vector $p_i$. Section 6 lists what is and is not shown.

\section{Setting and notation}
Categories $R=\{0,\dots,K-1\}$, $K=6$. Items $i=1,\dots,n$ (the page's $v\in V$; each item is $(x,y^\ast,\hat y)$ or a variant i--iv). Judge distribution $p_i=\Pj(\cdot\mid\cdot;\theta)\in\Delta^{K-1}$. Target $p^\ast_i=\Ps(\cdot\mid\cdot;h)$ for a human or persona $h$; with $k$ sparse ratings it is the empirical histogram, and $t_i$ is then the mean of the $k$ ratings (a noisy estimate of the true score). Means $m_i=\sum_r r\,p_i(r)$, $t_i=\sum_r r\,p^\ast_i(r)$; $\bar m,\bar t$ their averages over items; $\ttil_i=t_i-\bar t$.

\section{The equation per loss term}
{\small\renewcommand{\arraystretch}{1.15}
\begin{tabularx}{\textwidth}{@{}>{\raggedright\arraybackslash}p{1.75cm}>{\raggedright\arraybackslash}p{5.68cm}>{\raggedright\arraybackslash}X>{\raggedright\arraybackslash}X@{}}
\toprule
\textbf{Objective (to minimize)} & \textbf{Per-term equation} & \textbf{Reads of $p_i$} & \textbf{Coupled across items?}\\
\midrule
CE (page: sign dropped) &
$\begin{array}{@{}l@{}}\ell^{CE}_{i,r}=-p^\ast_i(r)\log p_i(r)\\ L_{CE}=\sum_i\sum_r\ell^{CE}_{i,r}\end{array}$ &
every coordinate $p_i(r)$ separately; unique population minimizer $p_i=p^\ast_i$ & no (sum of independent terms)\\
Corr &
$\begin{array}{@{}l@{}}L_{corr}=1-\rho,\quad \rho=\dfrac{N}{\sqrt{n}\,s_t\,\lVert m-\bar m\rVert}\\ N=\sum_i\sum_r r\,\ttil_i\,p_i(r),\quad s_t^2=\tfrac1n\sum_i\ttil_i^2\end{array}$ &
only $m_i=\sum_r r\,p_i(r)$; numerator term $(i,r)$ is $c_{i,r}\,p_i(r)$ with $c_{i,r}=r\,\ttil_i$ (linear); denominator through the spread of $m$ &
yes (via $\bar m$ and $\lVert m-\bar m\rVert$); equals $\frac1n\sum_i z_i\zeta_i$, the mean product of z-scores of the judge mean and the target mean\\
RPS &
$\displaystyle\sum_{k=0}^{K-2}\big(F_{p_i}(k)-F_{p^\ast_i}(k)\big)^2$ & whole CDF & no\\
TVD &
$\displaystyle\tfrac12\sum_r\lvert p_i(r)-p^\ast_i(r)\rvert$ & whole vector & no\\
\bottomrule
\end{tabularx}}

\section{Results (proofs by hand)}
\textbf{R1 (the formula in the box).}
\[
\sum_i(m_i-\bar m)(t_i-\bar t)=\sum_i\ttil_i m_i-\bar m\sum_i\ttil_i=\sum_i\ttil_i m_i
\qquad\text{because }\sum_i\ttil_i=0 .
\]
Insert $m_i=\sum_r r\,p_i(r)$. The denominator $\sqrt{\sum\ttil_i^2\sum(m_i-\bar m)^2}$ is a function of $m$ alone. \hfill$\blacksquare$

\textbf{R2 (Pearson sees only the means; level and scale are not identified).} By R1, $\rho=F(m_1,\dots,m_n)$. Two judges with the same $m_i$ for all $i$ have the same $\rho$ whatever their shapes; and
\[
\rho(a+bm)=\rho(m)\quad\text{for } b>0
\]
(numerator and denominator both scale by $b$, shifts cancel by centering). So the set of correlation-optimal judges is
\[
\{\,p:\ m_i=a+b\,t_i,\ b>0\,\}
\]
(when feasible inside $[0,K-1]$): $n(K-2)$ free shape dimensions plus $(a,b)$. A judge with the wrong level, the wrong scale and any spread is ``optimal''. \hfill$\blacksquare$

\textbf{R3 (gradients: Corr moves one direction per item, CE moves all of them).} With $p_i=\mathrm{softmax}(u_i)$:
\[
\frac{\partial L_{CE}}{\partial u_i(r)}=p_i(r)-p^\ast_i(r),\qquad
\frac{\partial\rho}{\partial u_i(r)}=p_i(r)\,(r-m_i)\,g_i,\qquad g_i=\frac{\zeta_i-\rho\,z_i}{n\,s_m},
\]
$z_i=(m_i-\bar m)/s_m$, $\zeta_i=\ttil_i/s_t$. Proof: $\partial\rho/\partial p_i(r)=r\,g_i$ (R1), and $\partial p_i(r')/\partial u_i(r)=p_i(r)(\delta_{rr'}-p_i(r'))$, so $\partial m_i/\partial u_i(r)=p_i(r)(r-m_i)$. \hfill$\blacksquare$ The correlation gradient on item $i$ always lies on the single direction $p_i\circ(r-m_i)$ (the score direction of the mean); the CE gradient ranges over the whole $(K-1)$-dimensional tangent space, so it is the one that can correct shape.

\textbf{R4 (exact decomposition of CE: the ``linear decomposition'').} Let $q_m$ be the tilted-uniform (max-entropy) distribution on $R$ with mean $m$, $G(P)=H(q_{m(P)})-H(P)\ge0$ (experiment 05). For any $P$ and any $q$,
\[
\mathrm{CE}(P,q)=\underbrace{H(P)}_{\text{irreducible noise}}+\underbrace{G(P)}_{\text{shape gap}}+\underbrace{\mathrm{KL}(q_{m(P)}\,\Vert\,q)}_{\text{level/scale/shape of }q}+\underbrace{\sum_r\big(P(r)-q_{m(P)}(r)\big)\big(-\log q(r)\big)}_{C(P,q)\ \text{(cross term)}}.
\]
Proof: $G+\mathrm{KL}(q_m\Vert q)=-H(P)+\mathrm{CE}(q_m,q)$, and $\mathrm{CE}(P,q)-\mathrm{CE}(q_m,q)=C$. \hfill$\blacksquare$ If $\log q$ is affine in $r$ (a judge that is a softmax of one affine score, a ``1-D'' judge), then $C=0$ and $\mathrm{KL}(q_{m(P)}\Vert q)$ is a function of the two means only: the term $G(P)$ is outside this log-affine family. Other scalar-to-distribution models can represent different shapes. A mean-based objective cannot separate forecasts that share each artifact’s mean. The paper's ``one-dimensional bottleneck'' is exactly $C=0$ plus $G>0$.

\textbf{R5 (RPS splits into a mean part and a shape part).} With $d_k=F_p(k)-F_q(k)$, $k=0,\dots,K-2$: $\sum_k d_k=m_q-m_p$ and
\[
\mathrm{RPS}(p,q)=\frac{(m_p-m_q)^2}{K-1}+\sum_k\big(d_k-\bar d\big)^2 .
\]
Proof: $\sum_k d_k^2=(\sum d_k)^2/(K-1)+\sum(d_k-\bar d)^2$ (variance identity), and $\sum_{k=0}^{K-2}F_p(k)=K-1-m_p$. \hfill$\blacksquare$ So RPS $\ge$ (mean error)$^2/(K-1)$, equality iff the CDF difference is constant. Likewise
\[
\lvert m_p-m_q\rvert\le(K-1)\,\mathrm{TVD}(p,q)
\]
(since $\sum_r r(p-q)(r)$, with $\sum(p-q)=0$, is at most $(K-1)$ times the positive part).

\textbf{R6 (propriety).} CE and RPS are strictly proper: the minimum over $p$ is attained only at $p=p^\ast$ (Gibbs; sum of squares of CDF gaps). TVD is zero only at $p=p^\ast$. $\rho$ is maximal on the whole set in R2, so it is not proper and cannot rank two judges whose means are affinely related to $t$. Concrete numbers (30 synthetic items, test against the truth, probabilities floored at $10^{-12}$, \texttt{math/results/verify\_math.json}): three judges all with $\rho=1.0000$:
\begin{center}
\begin{tabular}{@{}lrrr@{}}
\toprule
judge & CE & TVD & RPS\\
\midrule
truth & 1.269 & 0.000 & 0.000\\
max-entropy with the right mean & 1.698 & 0.369 & 0.082\\
two-point on $\{0,5\}$ with the right mean & 19.51 & 0.700 & 0.406\\
\bottomrule
\end{tabular}
\end{center}

\textbf{R7 (the literal reading of the page's last line).} If $\mathrm{Corr}(p^\theta,p^\ast)$ is taken as the correlation of the two histograms over $r$ (per item, 6 numbers each), it equals 1 for every
\[
p=\frac{1-b}{K}\mathbf 1+b\,p^\ast,\qquad b\in(0,1]:
\]
shrinkage towards uniform. Same conclusion (not proper), but this is not how correlation is used for judges; the box above (correlation of the judge's mean across items) is.

\section{Step-by-step derivations of R1--R6 (11-21-2026)}
Every symbol is as in the box on page 1: $p_i$ is the judge's distribution on item $i$, $m_i=\sum_r r\,p_i(r)$ its mean, $t_i$ the target mean, $\ttil_i=t_i-\bar t$, $n$ items, $R=\{0,\dots,K-1\}$.

\textbf{R1, step by step.}
\begin{enumerate}\setlength{\itemsep}{1pt}
\item Pearson correlation across items: $\rho=\dfrac{\sum_i(m_i-\bar m)(t_i-\bar t)}{\sqrt{\sum_i(m_i-\bar m)^2\sum_i(t_i-\bar t)^2}}$, and $t_i-\bar t=\ttil_i$.
\item Expand the numerator: $\sum_i(m_i-\bar m)\ttil_i=\sum_i m_i\ttil_i-\bar m\sum_i\ttil_i$.
\item $\sum_i\ttil_i=\sum_i t_i-n\bar t=0$, so the second term vanishes and the numerator is $\sum_i\ttil_i m_i$.
\item Insert $m_i=\sum_r r\,p_i(r)$: the numerator is $\sum_i\sum_r (r\,\ttil_i)\,p_i(r)$, linear in $p$ with fixed coefficients $c_{i,r}=r\,\ttil_i$.
\item In the denominator, $\sum_i\ttil_i^2$ is a constant of the target and $\sum_i(m_i-\bar m)^2$ depends on the judge only through $m$. Hence $\rho=F(m_1,\dots,m_n)$. \hfill$\blacksquare$
\end{enumerate}

\textbf{R2, step by step.}
\begin{enumerate}\setlength{\itemsep}{1pt}
\item Replace $m_i$ by $a+b\,m_i$ with $b>0$. The average becomes $a+b\,\bar m$, so every centered value becomes $b\,(m_i-\bar m)$ and $a$ cancels.
\item Numerator: $\sum_i\ttil_i\,b\,(m_i-\bar m)=b\sum_i\ttil_i m_i$ (R1).
\item Denominator: $\sqrt{\sum\ttil_i^2\cdot b^2\sum(m_i-\bar m)^2}=b\sqrt{\sum\ttil_i^2\sum(m_i-\bar m)^2}$, using $b>0$.
\item The factor $b$ cancels, so $\rho(a+b\,m)=\rho(m)$.
\item $\rho=1$ iff the centered vectors $(m_i-\bar m)$ and $(\ttil_i)$ are positively proportional (equality in Cauchy--Schwarz), i.e. $m_i-\bar m=b\,\ttil_i$ with $b>0$, i.e. $m_i=a+b\,t_i$ with $a=\bar m-b\,\bar t$.
\item Counting: $p_i$ has $K-1$ free parameters; fixing $m_i$ removes one, leaving $K-2$ per item, so $n(K-2)$ shape dimensions plus $(a,b)$ are free. \hfill$\blacksquare$
\end{enumerate}

\textbf{R3, step by step.}
\begin{enumerate}\setlength{\itemsep}{1pt}
\item Write $\rho=\dfrac{N}{\sqrt n\,s_t\,\lVert m-\bar m\rVert}$ with $N=\sum_i\ttil_i m_i$, $s_t^2=\frac1n\sum\ttil_i^2$, $s_m^2=\frac1n\sum(m_i-\bar m)^2$, so $\lVert m-\bar m\rVert=\sqrt n\,s_m$.
\item $\partial N/\partial m_j=\ttil_j$. Because $\sum_i(m_i-\bar m)=0$, the centering term drops and $\partial\lVert m-\bar m\rVert/\partial m_j=(m_j-\bar m)/\lVert m-\bar m\rVert$.
\item Quotient rule: $\dfrac{\partial\rho}{\partial m_j}=\dfrac{\ttil_j}{\sqrt n s_t\lVert m-\bar m\rVert}-\dfrac{N\,(m_j-\bar m)}{\sqrt n s_t\lVert m-\bar m\rVert^3}=\dfrac{\zeta_j-\rho\,z_j}{n\,s_m}=g_j$, with $\zeta_j=\ttil_j/s_t$, $z_j=(m_j-\bar m)/s_m$.
\item $\partial m_i/\partial p_i(r)=r$, so $\partial\rho/\partial p_i(r)=r\,g_i$.
\item Softmax $p_i(r')=e^{u_i(r')}/\sum_s e^{u_i(s)}$ gives $\partial p_i(r')/\partial u_i(r)=p_i(r')(\delta_{rr'}-p_i(r))$. Hence $\dfrac{\partial m_i}{\partial u_i(r)}=\sum_{r'}r'p_i(r')(\delta_{rr'}-p_i(r))=p_i(r)\,(r-m_i)$.
\item Chain rule: $\partial\rho/\partial u_i(r)=p_i(r)(r-m_i)\,g_i$: the vector $p_i\circ(r-m_i)$ times the single number $g_i$.
\item CE: $L_{CE}=-\sum_i\sum_r p^\ast_i(r)\log p_i(r)$, so $\dfrac{\partial L_{CE}}{\partial u_i(r)}=-\sum_{r'}\dfrac{p^\ast_i(r')}{p_i(r')}\,p_i(r')(\delta_{rr'}-p_i(r))=-p^\ast_i(r)+p_i(r)\sum_{r'}p^\ast_i(r')=p_i(r)-p^\ast_i(r)$. \hfill$\blacksquare$
\end{enumerate}

\textbf{R4, step by step.}
\begin{enumerate}\setlength{\itemsep}{1pt}
\item Let $q_m$ be the maximum-entropy distribution on $R$ with mean $m$ (an exponential tilt of the uniform: $q_m(r)\propto e^{\lambda r}$), and $m(P)$ the mean of $P$. Define $G(P)=H(q_{m(P)})-H(P)$; it is $\ge0$ because $q_{m(P)}$ maximizes entropy among all distributions with mean $m(P)$, and $P$ is one of them.
\item Add and subtract: $\mathrm{CE}(P,q)=\mathrm{CE}(q_{m(P)},q)+C(P,q)$ with \[C(P,q)=\sum_r\big(P(r)-q_{m(P)}(r)\big)(-\log q(r)).\]
\item $\mathrm{CE}(q_m,q)=H(q_m)+\mathrm{KL}(q_m\Vert q)$ (the identity $\mathrm{CE}=H+\mathrm{KL}$), and $H(q_m)=H(P)+G(P)$.
\item Together: $\mathrm{CE}(P,q)=H(P)+G(P)+\mathrm{KL}(q_{m(P)}\Vert q)+C(P,q)$.
\item If $-\log q(r)=\alpha+\beta r$ (log-affine, a ``1-D'' judge): $C=\alpha\sum_r(P-q_m)(r)+\beta\sum_r r\,(P-q_m)(r)=\alpha\cdot0+\beta\,(m(P)-m(P))=0$, since $P$ and $q_{m(P)}$ both sum to 1 and both have mean $m(P)$.
\item Then $\mathrm{CE}(P,q)=H(P)+G(P)+\mathrm{KL}(q_{m(P)}\Vert q)$: the only part $q$ controls is the KL term, a function of the two means. $G(P)$ is the same for every 1-D judge. \hfill$\blacksquare$
\end{enumerate}

\textbf{R5, step by step.}
\begin{enumerate}\setlength{\itemsep}{1pt}
\item $\sum_{k=0}^{K-2}F_p(k)=\sum_{k=0}^{K-2}\sum_{r\le k}p(r)=\sum_r p(r)\,\#\{k:\ r\le k\le K-2\}=\sum_r p(r)(K-1-r)=K-1-m_p$ (for $r=K-1$ the count is $0=K-1-r$).
\item Hence $\sum_k d_k=\sum_kF_p(k)-\sum_kF_q(k)=m_q-m_p$.
\item Variance identity over the $K-1$ values $d_k$ with mean $\bar d=\frac{1}{K-1}\sum_kd_k$: $\sum_kd_k^2=(K-1)\bar d^{\,2}+\sum_k(d_k-\bar d)^2=\dfrac{(\sum_kd_k)^2}{K-1}+\sum_k(d_k-\bar d)^2$.
\item With $\mathrm{RPS}=\sum_kd_k^2$ and step 2: $\mathrm{RPS}(p,q)=\dfrac{(m_p-m_q)^2}{K-1}+\sum_k(d_k-\bar d)^2$. Equality with the first term iff all $d_k$ are equal.
\item TVD bound: $m_p-m_q=\sum_r r\,(p-q)(r)=\sum_r\big(r-\tfrac{K-1}2\big)(p-q)(r)$ because $\sum_r(p-q)(r)=0$. Since $\lvert r-\tfrac{K-1}2\rvert\le\tfrac{K-1}2$, $\lvert m_p-m_q\rvert\le\tfrac{K-1}2\sum_r\lvert p-q\rvert=(K-1)\,\mathrm{TVD}(p,q)$. \hfill$\blacksquare$
\end{enumerate}

\textbf{R6, step by step.}
\begin{enumerate}\setlength{\itemsep}{1pt}
\item CE: $\mathrm{CE}(p^\ast,p)-\mathrm{CE}(p^\ast,p^\ast)=\mathrm{KL}(p^\ast\Vert p)\ge0$, with equality iff $p=p^\ast$ (Gibbs' inequality). So CE is minimized only at $p^\ast$.
\item RPS $=\sum_kd_k^2=0$ iff every CDF gap is $0$ iff $F_p=F_{p^\ast}$ iff $p=p^\ast$.
\item TVD $=\frac12\sum_r\lvert p(r)-p^\ast(r)\rvert=0$ iff $p=p^\ast$.
\item $\rho\le1$, and by R2 it equals 1 on the whole set $\{p:\ m_i=a+b\,t_i,\ b>0\}$, which contains many $p\ne p^\ast$. So $\rho$ is not uniquely maximized at $p^\ast$, and it cannot rank two judges whose means are affinely related to $t$. \hfill$\blacksquare$
\end{enumerate}

\section{Evaluation, as the page says it is ``wrong''}
Using correlation of the judge's mean with human mean ratings as the evaluation metric has the R2 blindness: level, scale and spread are invisible, and three judges with $\rho=1$ (R6) differ by 18 nats of CE. The metrics that read the whole distribution are test CE, test TVD and test RPS (all lower is better), with Pearson and Spearman as sanity columns; Sanmi's point 4 (distributional accuracy when the output is read as a rating or distribution; rank agreement for ranking uses) matches. Caveat from the page: judges selected on Corr will tend to win on Corr metrics, so a Corr win on a Corr metric says little; the sweep reports both.

\section{Checks, and what the result does and does not show}
\small\sloppy\renewcommand{\arraystretch}{1.3}
\begin{tabularx}{\textwidth}{@{}>{\raggedright\arraybackslash}p{2.6cm}>{\raggedright\arraybackslash}p{3.6cm}>{\raggedright\arraybackslash}X@{}}
\toprule
\textbf{Check} & \textbf{Where} & \textbf{Result}\\
\midrule
By hand & Section 3 & R1--R7\\
sympy exact ($n=3$, $K=4$ / $K=4$ / $K=6$) & \texttt{math/\allowbreak verify\_math.py} S1--S4 & numerator linear in $p$ (exact); $\partial\rho/\partial p_i(r)=r g_i$ (exact derivative vs formula, diff 0 at a rational point); RPS split and $\sum d_k=m_q-m_p$ (exact); vector-Pearson identity (exact)\\
numerical, $K=6$, random & \texttt{math/\allowbreak verify\_math.py} N1--N6 & shape changes leave $\rho$ unchanged (max diff $3\mathrm{e}{-16}$, 200 draws); affine invariance ($3\mathrm{e}{-16}$); logit gradients of $\rho$ and CE match finite differences ($2\mathrm{e}{-10}$, $<10^{-6}$); CE decomposition R4 ($8.9\mathrm{e}{-16}$, 300 draws); $C=0$ for 1-D judges ($10^{-12}$); RPS split ($1.8\mathrm{e}{-15}$, 500 pairs); $\lvert\Delta m\rvert\le(K-1)$TVD\\
Lean 4.22 + Mathlib & \texttt{math/\allowbreak CEvsCorr.lean}, log in \texttt{math/\allowbreak results/\allowbreak lean\_check.log} & proved, no \texttt{sorry}: mean is linear in $p$; Pearson numerator needs no centering of the judge means; $\rho$ computed from the means is equal for judges with equal means; a concrete pair (uniform on $\{1,2,3\}$, uniform on $\{0,2,4\}$) with the same mean, different shape, RPS $2/9>0$; CDF-difference sum equals mean difference; RPS split; RPS $\ge$ (mean error)$^2/5$\\
\bottomrule
\end{tabularx}

\normalsize\textbf{Shows:} with the human distribution known, Pearson-of-means is a function of the judge's means only (R1--R3), is blind to level, scale and shape, is not proper (R6), and the CE loss decomposes exactly into noise, an unreachable shape gap, and a level/scale/shape-of-$q$ term (R4); RPS has the same mean-plus-shape split (R5).

\textbf{Does not show:} (1) That a CE-trained judge generalizes better with few ratings. That is empirical (exp 02, 08, 11, 16, 17 and the sweep in this folder). (2) That correlation \emph{training} of a flexible judge fails: at the population level it is not proper, but with finite data it can act as a regularizer; the sweep tests when. (3) Anything about the real rating noise: $t_i$ here is the target mean; with 2--3 raters it is estimated (Sanmi point 4). (4) The Lean file proves the building blocks; the full theorem ``$\rho(a+bm)=\rho(m)$'' and the R4/R3 identities are checked with sympy and numbers, not in Lean. (5) Spearman and Kendall depend on the ranks of $m$, so they share the ``only the means'' property but are not covered by the formula in the box.

\textbf{TLDR-end:} [exp 18 math] Pearson between rating means depends only on the first conditional moment, up to a positive affine change; CE splits exactly into noise, an unreachable shape gap and a calibration term; checked by hand, sympy, numbers and Lean.

\end{document}
